\paragraph{Exact output.}\label{sec:cr-exact}

\begin{lemma}[Uniform height over residual labels]\label{lem:cr-height}
Let $F$ be a rational quadratic with continuous core $[0,1]^k$ and residual
satisfying (R1), with integer residual bounds rounded inward. Let
$\Lambda_c$ be the least common denominator of all coefficients of $F$,
$\Lambda_e$ that of all residual bounds, $T_0=\Lambda_c\Lambda_e^2$,
$H_0=\max\{1,\max_{i,j\in\mathcal C}|T_0\,\partial_{ij}F|\}$ and
\[
D_0=T_0\bigl(k!\,H_0^k\bigr)^2 .
\]
For every residual endpoint label $z$, the value
$\Phi(z)=\min_{v\in[0,1]^k}F(v,z)$ is a rational with reduced denominator
at most $D_0$, and so is $\OPT$.
\end{lemma}
\begin{proof}
Fix $z$ and put $\phi(v)=F(v,z)$. Every monomial of $T_0\phi$ has an
integer coefficient: the coefficients of $v_iv_j$ and $v_i^2$ are
coefficients of $F$; the coefficient of $v_i$ is $b_i+\sum_jc_{ij}z_j$;
and the constant collects $c_0$, $a_jz_j^2$, $b_jz_j$ and $c_{ij}z_iz_j$
for residual $i,j$. Each has denominator dividing $\Lambda_c\Lambda_e^2$.
The matrix $A=T_0\nabla^2_{\mathcal C\mathcal C}F$ is integral with entries
bounded by $H_0$ in absolute value.

Among minimizers of $\phi$ on $[0,1]^k$ choose one, $v$, with the fewest
coordinates in $(0,1)$; let $S$ be that set. The other coordinates are
$0$ or $1$. If $S=\emptyset$, then $T_0\phi(v)\in\Z$. Otherwise
$\nabla_S\phi(v)=0$ and $A_{SS}\succeq0$ by first- and second-order
optimality on the face. If $A_{SS}$ were singular, a nonzero $w$
supported on $S$ with $A_{SS}w_S=0$ would give $\phi(v+tw)=\phi(v)$ for all
$t$, and moving $t$ until a coordinate of $S$ reaches $0$ or $1$ would
produce a minimizer with fewer interior coordinates. Hence
$\delta=|\det A_{SS}|$ is a positive integer, and by the Leibniz expansion
$\delta\le|S|!\,H_0^{|S|}\le k!\,H_0^k$. The stationarity system
$A_{SS}v_S=-T_0\nabla_S\phi(0_S,v_{\mathcal C\setminus S})$ has an integral
right-hand side, so $\delta v$ is integral by Cramer's rule. Evaluating
the integer-coefficient quadratic $T_0\phi$ at $v=(\delta v)/\delta$ gives
an element of $\delta^{-2}\Z$. Thus $\Phi(z)\in(T_0\delta^2)^{-1}\Z$.
By Lemma~\ref{lem:cr-endpoint}, some global optimizer has an endpoint
label $z^*$, and $\OPT=\Phi(z^*)$.
\end{proof}

\begin{remark}
Any common multiple $S_0$ of all denominators can replace $\Lambda_c$ and
$\Lambda_e$, for example $T_0=S_0^3$. A cube, or the mixed form
$\Lambda_c\Lambda_e^2$, is needed in general: the term $\frac12x_ix_j$ at
residual endpoints $x_i=x_j=\frac12$ equals $\frac18$, while the least
common denominator of the data is $2$.
\end{remark}

\begin{theorem}[Exact output with a cut residual]\label{thm:cr-exact}
Assume~\eqref{eq:cr-class} with continuous core $[0,1]^k$ and rational data.
Run the core search until $U-\min\{\lambda_j,U\}<D_0^{-2}$, which happens
by the first level with $e_j<D_0^{-2}$. Let $z_U$ be the residual label
of the incumbent completion. Minimize $\phi(v)=F(v,z_U)$ exactly over
$[0,1]^k$ by enumerating the $3^k$ faces of the cube, solving the
stationarity system on each face whose free Hessian block is nonsingular,
and keeping feasible solutions. The best candidate $\hat v$ satisfies
$F(\hat v,z_U)=\OPT$, so $(\hat v,z_U)$ is a global optimizer. A checker
needs only the search trace, $D_0$, feasibility and exact value of the
output, its reduced denominator ($\le D_0$), and the inequality
$F(\hat v,z_U)-\min\{\lambda_j,U\}<D_0^{-2}$.
If in addition $V$ has core growth with constant $g$ as in
Proposition~\ref{prop:cr-growth}, the total bit complexity is
$f(k,\kappa)(I+1)^{C_1}$ with an absolute constant $C_1$, and $g$ is not
an input.
\end{theorem}
\begin{proof}
The proof of Lemma~\ref{lem:cr-height} shows that a minimizer of $\phi$
lies on some face with a nonsingular free block, so the enumeration
contains it; all candidates are feasible, so the best one attains
$\Phi(z_U)$. If $(v_U,z_U)$ is the incumbent completion, then
$\OPT\le\Phi(z_U)\le F(v_U,z_U)=U$ and, by Theorem~\ref{thm:cr-search}(d),
$\min\{\lambda_j,U\}\le\OPT$. So $0\le\Phi(z_U)-\OPT<D_0^{-2}$. Both
numbers have reduced denominators at most $D_0$ by
Lemma~\ref{lem:cr-height}, and distinct such rationals differ by at least
$D_0^{-2}$; hence they are equal. The checker's test proves optimality of
any feasible point with these properties by the same separation argument.

For the complexity, $\log T_0$ and $\log H_0$ are $O(I)$ and $k\le I$, so
$\log D_0=O(I^2)$ and the number of levels is
$J=O(\log D_0+\log(kL)+1)=\poly(I)$. Under core growth,
Proposition~\ref{prop:cr-growth} bounds the queries by
$2^k+8^kJ(1+\sqrt{k\kappa})^k$. Each query is a cut problem whose data have
bit length polynomial in $I+J$, hence $\poly(I)$ by
Theorem~\ref{thm:cr-oracle}. Face enumeration costs $3^k\poly(I)$. All
factors depending only on $k$ and $\kappa$ form $f$.
\end{proof}

Without core growth the exact mode still terminates, but its number of
cells can be $2^{kJ}$ with $J=\Theta(\log D_0)$. The stopping depth must
not be combined with Theorem~\ref{thm:cr-smoothed} by choosing
$M\ge2^{J}$: sampled coefficients enter $D_0$, and for a fixed noise scale
the atom denominators grow with $M$, so this condition generally cannot
hold.

